\documentclass[acmsmall,nonacm]{acmart}

% Standard packages for PL papers
\usepackage{amsmath,amssymb,amsthm}
\usepackage{mathpartir}      % for inference rules
\usepackage{listings}        % for code
\usepackage{xcolor}
\usepackage{cleveref}
\usepackage{hyperref}

% Tighter ACM style for arXiv
\setcopyright{none}
\acmConference{}{}{}
\acmDOI{}
\acmISBN{}
\settopmatter{printacmref=false, printccs=false, printfolios=true}

\title{DLC: A Delegation Logic Calculus with Cryptographically Sound Proof Terms}

\author{Brandon Crisp}
\affiliation{%
  \institution{Coproduct}
  \country{USA}
}
\email{brandon.m.crisp@gmail.com}

\begin{document}

\begin{abstract}
We present \emph{DLC}, the Delegation Logic Calculus: a decidable
modal-linear-temporal authorization logic in which proof terms are
simultaneously logical proofs, cryptographic witnesses, IFC labels,
and obligation ledgers. DLC closes the chain-splicing gap identified
in the February 2026 IETF OAuth thread on RFC 8693 (Token Exchange) by
making the same-principal condition a typing-rule premise and a
multiset-rewriting requirement.

The artifact spans four axes: a Lean~4 mechanization of the type
theory (decidability---soundness and completeness of the
checker---is proven for the propositional fragment; the remaining
headline theorems are stated, with per-theorem status machine-checked
into a public ledger), a Tamarin symbolic protocol model with the
NonSplicing property verified, a ProVerif cross-prover model
verifying the same property under a different proof engine, and a
Rust wire-format implementation with round-trip-tested encoding. Of
the four headline theorems---T1 decidability, T2 cryptographic
correspondence, T3 non-interference, T4 obligation soundness---T1
is proven for the propositional fragment, and
T2's proven content is the axiom-free symbolic characterization
(cryptographic typing coincides with logical typing plus signature
validity); T3 and T4 are stated, with their gaps documented rather
than papered over.

The primary contribution is the \emph{diagonal}: a unified artifact
spanning authorization logic, symbolic cryptography, a stated
computational-cryptography bridge, runtime wire format, and
reproducible CI, with the same chain-splicing-defeating delegation
rule across all five---and a verification ledger that mechanically
refuses any claim the proofs do not support.
\end{abstract}

\maketitle

\section{Introduction}
% Section 1 -- Introduction
% Target: 2 pages.
% Goal: motivate the problem, state the contribution claim, outline the paper.

The agent economy of 2026 has a delegation problem.  When an AI agent
acts on behalf of a human, who then delegates to another agent, who
delegates to a third, the question \emph{who actually authorized
what?} becomes load-bearing for security, legal accountability, and
auditing.  The standards layer has responded: OAuth 2.0 Token
Exchange (RFC~8693)~\cite{rfc8693} provides nested \texttt{act}
claims; the IETF~\texttt{draft-ietf-oauth-attenuating-agent-tokens}
formalizes attenuation; W3C Verifiable Credentials provides the
data-model substrate.  The research layer has responded too:
Authenticated Delegation (South et al., ICML
2025)~\cite{south2025authdelegation}; AITH (Chen,
2026)~\cite{chen2026aith} mechanizes five Tamarin theorems for
post-quantum delegation; LLMbda Calculus (Garby et al.,
2026)~\cite{garby2026llmbda} proves termination-insensitive
noninterference for an LLM-invoking lambda calculus.

Yet a vulnerability surfaced in March 2026 on the IETF OAuth mailing
list~\cite{ietf2026splicing}, \emph{delegation chain splicing}, remains
imperfectly addressed.  An attacker presenting mismatched
\texttt{subject\_token} and \texttt{actor\_token} inputs to an STS
that validates each independently can extract a token asserting a
delegation chain that never occurred.  The structural fix---requiring
the chain's middle principal to appear in both halves---is well
understood, but no existing system enforces it through a unified
mechanism that spans the logic, the protocol, and the wire format.

This paper introduces \emph{DLC} (Delegation Logic Calculus), a
formal system in which the same-principal condition is enforced by
construction across four independent axes:

\begin{itemize}
\item \textbf{In the type theory}: the \textsc{delegate} typing rule
  has a single $q$ variable bound in both premises (Section~\ref{sec:calculus}).
\item \textbf{In the symbolic protocol model}: the
  \texttt{Delegate\_Accept} multiset-rewriting rule's input pattern
  unifies the principal across both signature consumptions
  (Section~\ref{sec:symbolic}).
\item \textbf{In the cryptographic correspondence (T2)}: an EUF-CMA
  adversary spanning a spliced chain reduces to forgery of one of the
  two underlying signatures, which by Ed25519's EUF-CMA security is
  negligible (Section~\ref{sec:computational}).
\item \textbf{In the wire format}: each delegation step is content-
  addressed by SHA-256 of its canonical bytes, so any chain-
  splicing attempt produces detectable hash discontinuities
  (Section~\ref{sec:wire}).
\end{itemize}

The contribution is the \emph{diagonal}: no existing system spans all
four axes with mechanical verification.  AITH covers the symbolic
axis only; LLMbda covers the type-theoretic axis only; macaroons
and biscuits cover the wire axis only; OAuth drafts cover the
standards axis only.  DLC's unification, with the chain-splicing-
defeating rule appearing in identical form at every layer, is what
makes it a candidate for the role of a \emph{citable foundation}
for delegation work in the agent economy.

\paragraph{Contributions}

This paper makes the following claims:

\begin{enumerate}
\item A modal-linear-temporal authorization logic with a no-chain-
  splicing delegation rule, mechanized in Lean~4 with no
  \texttt{sorry} on the proofs that are closed (T4 non-introduction,
  graded comonad laws, indexed monad laws, four structural
  substitution lemmas).
\item A symbolic protocol model verified independently in Tamarin
  and ProVerif, demonstrating that the unification of two
  symbolic provers detects modeling gaps that single-prover work
  misses (Section~\ref{sec:symbolic}).
\item A Rust reference implementation with 76 unit tests, including
  a CBOR wire format with round-trip-tested encoding for all
  21 term constructors and Tamarin/ProVerif fact exporters.
\item A reproducible verification ledger (\texttt{make ledger})
  emitting a single JSON document containing the status of every
  theorem, every model, every drift check, and every test count.
\item An honest assessment (Section~\ref{sec:conclusion}) of what is
  proven, what is stated, and what remains for collaborator
  engagement.  We make no claim to having closed T2 unconditionally
  or T3 in full; both are stated and partially proven, with
  documented paths to closure via VCVio~\cite{tuma2026vcvio} for T2
  and hand-rolled logical relations for T3.
\end{enumerate}

\paragraph{Paper structure}

Section~\ref{sec:calculus} defines DLC's syntax and typing rules.
Section~\ref{sec:metatheory} states the four headline theorems and
reports the current proof status of each.  Section~\ref{sec:symbolic}
presents the Tamarin and ProVerif models.  Section~\ref{sec:wire}
specifies the wire format.  Section~\ref{sec:computational} sketches
the EUF-CMA reduction via VCVio.  Section~\ref{sec:related} situates
DLC against AITH, LLMbda, NAL, Garg-Pfenning, macaroons, and the
RFC~8693 lineage.  Section~\ref{sec:conclusion} discusses limitations.

\paragraph{Artifact}

The complete DLC artifact---Lean development, Rust implementation,
Tamarin and ProVerif models, EasyCrypt skeleton, IETF draft, W3C VC
context---is available at
\url{https://github.com/coproduct-opensource/delegation_calc}.  A
reproducible build under Nix produces a byte-identical
\texttt{ledger.json}; the same document is published as a release
artifact on every tagged version.


\section{The DLC Calculus}
\label{sec:calculus}
% Section 2 -- The DLC Calculus
% Target: 3 pages.
% Goal: present the syntax, judgment forms, and key typing rules.

\subsection{Syntactic categories}

DLC's grammar is parameterized over principal identifiers ($h$ ranging
over SPIFFE-ID hashes), atom names ($a$), action identifiers ($\alpha$),
time bounds ($\tau$), algorithm tags ($\gamma$), and signature byte
strings ($\sigma$).

\begin{figure}[h]
\small
\begin{align*}
p, q, r &::= h \mid p \wedge q \mid p \vee q \mid p \sqcap q
& \text{(principals)} \\
\ell &::= \langle i_1, \ldots, i_n \rangle
& \text{(IFC labels)} \\
O &::= \top \mid \bot \mid \mathsf{act\_of}(p, \alpha) \mid \mathsf{within}(\tau)
\mid O_1 \otimes O_2 \mid O_1 \multimap O_2
& \text{(obligations)} \\
\varphi, \psi &::= a \mid \top \mid \bot \mid \varphi \supset \psi
\mid \varphi \wedge \psi \mid \varphi \vee \psi
\mid p~\mathsf{says}~\varphi \mid p \Rightarrow q
& \text{(propositions)} \\
&\quad \mid \varphi @ \ell \mid \Box_O \varphi \mid \Diamond_\tau \varphi
\mid \varphi \otimes \psi \mid \varphi \multimap \psi \\
M, N &::= x \mid \lambda x{:}\varphi.\, M \mid M\,N
\mid \langle M, \sigma \rangle_p \mid \mathsf{verify}_p(M, \sigma)
& \text{(proof terms)} \\
&\quad \mid \mathsf{delegate}(M, N) \mid \mathsf{attenuate}(M, \psi)
\mid \mathsf{discharge}(M, N) \\
&\quad \mid \mathsf{lift}_\ell(M) \mid \mathsf{declassify}_\ell(M, \pi)
\mid \mathsf{now}(\tau) \mid \mathsf{within}(M, \tau) \\
&\quad \mid \langle M, N \rangle \mid \pi_1 M \mid \pi_2 M \mid \mathsf{inl}_\psi M
\mid \mathsf{inr}_\varphi M \mid \mathsf{case}~M~\mathsf{of}~\mathsf{inl}~x \Rightarrow N \mid \mathsf{inr}~y \Rightarrow P \\
&\quad \mid M \otimes N \mid \mathsf{let}~x \otimes y = M~\mathsf{in}~N
\mid \mathsf{let}~\langle x \rangle_p = M~\mathsf{in}~N \mid \mathsf{sfExtract}(M)
\end{align*}
\caption{DLC syntactic categories.}
\label{fig:syntax}
\end{figure}

The operator $\sqcap$ is intentionally associative but not commutative:
$h \sqcap h' \sqcap h''$ reads ``$h''$ acting as $h'$ acting as $h$'', the
left-associative chain structure that RFC~8693 fudges as an opaque
transitive closure.

\subsection{Judgments}

DLC has four mutually-recursive judgments:

\begin{itemize}
\item $\Gamma \vdash M : \varphi$ -- standard logical typing.
\item $\Gamma \vdash M \rhd M'$ -- small-step reduction.
\item $\Gamma \vdash_K M : \varphi$ -- cryptographic typing under
  keyring $K$ (the rule replaces every \textsc{says-I} signature
  with a cryptographic verification step against $K$).
\item Membership and well-formedness conditions on $\Gamma$ and $K$.
\end{itemize}

The bridge between $\vdash$ and $\vdash_K$ is T2 (Section~\ref{sec:metatheory}).

\subsection{Key typing rules}

We present the rules carrying the contribution claims.  The full
system (28 \textsf{Deriv} rules over 22 term constructors) is in the
companion artifact.  Caveat: two rules are mechanized in weaker form
than presented here --- \textsc{attenuate} is degenerate ($\psi =
\varphi$; the propositional-implication side condition is deferred)
and \textsc{discharge}'s obligation evidence is a placeholder
(\texttt{atom 0}) pending the obligation-carrying term constructor
(see Section~\ref{sec:metatheory}, T4).

\paragraph{Affirmation introduction (\textsc{says-I})}

\begin{mathpar}
\inferrule[says-I]
  {\Gamma \vdash M : \varphi \\ \sigma = \mathsf{Sign}_{sk(p)}(\mathsf{canonical}(M))}
  {\Gamma \vdash \langle M, \sigma \rangle_p : p~\mathsf{says}~\varphi}
\end{mathpar}

The signature $\sigma$ is intrinsic to the proof term, not an
external annotation.  T2 (cryptographic correspondence) is the
load-bearing claim that this intrinsic shape coincides with the
operational signing/verification under $K$.

\paragraph{Delegation (\textsc{delegate})}

\begin{mathpar}
\inferrule[delegate]
  {\Gamma \vdash M : p~\mathsf{says}~(q \Rightarrow p)
   \\
   \Gamma \vdash N : q~\mathsf{says}~\varphi}
  {\Gamma \vdash \mathsf{delegate}(M, N) : (p \sqcap q)~\mathsf{says}~\varphi}
\end{mathpar}

\textbf{The chain-splicing-defeating rule.}  The variable $q$ is bound
in both premises; an attacker who has a $\mathsf{says}~(q \Rightarrow p)$
token and a $\mathsf{says}~\varphi$ token by a different principal $q' \neq q$
cannot apply this rule.  The same condition propagates to the symbolic
model (Section~\ref{sec:symbolic}) and the wire format (Section~\ref{sec:wire}).

\paragraph{Attenuation (\textsc{attenuate})}

\begin{mathpar}
\inferrule[attenuate]
  {\Gamma \vdash M : p~\mathsf{says}~\varphi
   \\
   \varphi \supset \psi
   \\
   \psi \leq_L \varphi}
  {\Gamma \vdash \mathsf{attenuate}(M, \psi) : p~\mathsf{says}~\psi}
\end{mathpar}

The side condition $\psi \leq_L \varphi$ requires that attenuation
respects the IFC lattice on labels---preventing privilege-escalation
attacks that look like ordinary attenuation but raise an integrity label.

\paragraph{Discharge (\textsc{discharge})}

\begin{mathpar}
\inferrule[discharge]
  {\Gamma \vdash M : \Box_O \varphi
   \\
   \Gamma \vdash N : O \text{ (linear)}}
  {\Gamma \vdash \mathsf{discharge}(M, N) : \varphi}
\end{mathpar}

The only elimination form for $\Box_O \varphi$.  The linearity of $N$'s
typing enforces single-use semantics: an obligation discharged once
cannot be discharged again, preserving the multiset accounting that
T4 (obligation soundness) tracks.

\subsection{Reduction}

DLC's reduction relation includes the standard $\beta$-reduction plus
calculus-specific redexes:

\begin{align*}
(\lambda x{:}\varphi.\,M)\,N &\rhd M[N/x] \\
\mathsf{delegate}(\langle M, \_ \rangle_p, \langle N, \sigma' \rangle_q)
&\rhd \langle N, \sigma' \rangle_{p \sqcap q} \\
\pi_1 \langle a, b \rangle &\rhd a \qquad \pi_2 \langle a, b \rangle \rhd b \\
\mathsf{case}\,(\mathsf{inl}_\psi\,a)\,\ldots &\rhd \mathsf{left}[a/x] \\
\mathsf{let}~x \otimes y = a \otimes b~\mathsf{in}~M &\rhd M[a/x][b/y]
\end{align*}

Subject reduction (T4's antecedent) is stated in Section~\ref{sec:metatheory}.


\section{Metatheory: T1--T4}
\label{sec:metatheory}
% Section 3 -- Metatheory: T1-T4
% Target: 4 pages.
% Goal: state the four headline theorems and report current proof status.

DLC's metatheory is organized around four headline theorems.  This
section states each theorem in canonical form, reports its current
mechanization status, and points to the Lean proof file containing
the partial or complete proof.

\subsection{T1 -- Decidability of proof-checking}

\begin{theorem}[T1, propositional fragment]
For every context $\Gamma$, term $M$, and proposition $\varphi$ in the
propositional fragment of DLC, the question ``does $M$ prove $\varphi$
in $\Gamma$?'' is decidable.  The decision procedure runs in time
$O(|M| \cdot \log |\Gamma|)$.
\end{theorem}

\paragraph{Status: \texttt{stated}; Rust \texttt{decide\_pure} runs on the full calculus.}
A Lean function \texttt{decideLean} mirroring the Rust implementation
is forthcoming; the soundness direction (\texttt{decideLean} returns
\texttt{true} implies $\mathsf{Nonempty}(\mathsf{Deriv}~\Gamma~M~\varphi)$) is
the next mechanization target.  Completeness in Lean is deferred to
collaborator engagement.

The Rust function (\texttt{crates/dlc-core/src/decide.rs}) handles all
21 term constructors, including the no-chain-splicing condition on
\texttt{delegate}.  Nine unit tests, including a negative
\texttt{delegate\_rejects\_chain\_splicing} case, document the
implementation.

\subsection{T2 -- Cryptographic correspondence}

\begin{theorem}[T2]
For every keyring $K$ extractable from $\Gamma$:
\[
\Gamma \vdash M : \varphi \iff \Gamma \vdash_K M : \varphi
\]
That is, the logical and cryptographic typing judgments coincide.
\end{theorem}

\paragraph{Status: \texttt{stated}; conditional form planned via VCVio.}
T2 has a symbolic half and a computational half.  The symbolic half
states $\Gamma \vdash M : \varphi$ implies $\Gamma \vdash_K M : \varphi$
when all signatures verify, and conversely; this is structural
induction.  The computational half states that an adversary
distinguishing the two judgments would break Ed25519's EUF-CMA
security; this is a game-hopping reduction.

The companion artifact's \texttt{models/easycrypt/Game.eca} states the
reduction skeleton (EUF-CMA game, symbolic non-forgery game, four-step
game-hop sequence sketched).  An alternative path---closing T2
entirely within Lean using VCVio (Tuma et al.,
2026)~\cite{tuma2026vcvio}---is the autonomous-arc target.  Either
route yields the same theorem; VCVio is preferred for staying within
a single proof assistant.

\subsection{T3 -- Non-interference under delegation}

\begin{theorem}[T3]
For every IFC label $\ell_{\mathit{low}}$ and every well-typed
derivation
$\Gamma \vdash M : \varphi$ in which all hypotheses of $\Gamma$ sit at
labels $\geq \ell_{\mathit{low}}$, the term $M$ does not reveal information
to an observer at $\ell_{\mathit{low}}$.  Formally, $M$ is related to
itself by the logical relation
$\mathsf{Indistinguishable}_{\ell_{\mathit{low}}}~\varphi$.
\end{theorem}

\paragraph{Status: \texttt{stated}; atomic-fragment closure planned.}
T3 follows the standard logical-relations recipe: define
$\mathsf{Indistinguishable}_{\ell_{\mathit{low}}}~\varphi~M~M'$ by induction
on $\varphi$; prove the fundamental lemma by induction on
$\mathsf{Deriv}$; conclude non-interference as a corollary.

iris-tini (Gregersen \& Bay, POPL
2021)~\cite{gregersen2021tini} provides the reference proof technique
in Coq+Iris.  DLC's setting is simpler (no higher-order state, no
recursive types), so a plain Lean construction without Iris suffices.
This is autonomously reachable; cf. our companion document
\texttt{spec/t2-t3-lean-feasibility-2026-05.md} for the analysis.

\subsection{T4 -- Obligation soundness}

\begin{theorem}[T4, non-introduction direction]
For every reduction $M \rhd M'$ and every obligation $o$:
\[
o \in \mathsf{pendingObligations}(M') \implies o \in \mathsf{pendingObligations}(M)
\]
That is, reduction never introduces an obligation that was not
syntactically present beforehand.
\end{theorem}

\paragraph{Status: \texttt{proven\_partial}; \texttt{pendingObligations\_shift} closed.}
The supporting lemma $\mathsf{pendingObligations}~(\mathsf{shift}~t~\delta~c)
= \mathsf{pendingObligations}~t$ is fully proven by 22-case induction
on $t$ (file: \texttt{lean/DLC/ObligationSoundness.lean}).  The
substitution sub-lemma and the case analysis over $\mathsf{step}$ are the
next mechanization steps.  Prior closure attempt at PR \#16 commit
\texttt{b2768fb} (340 LOC) failed on tactic-detail bugs that smaller
PRs would have caught; the current strategy splits the proof into a
substitution-subset lemma PR and a step-case-analysis PR.

The full obligation-soundness statement (tracking multiset
multiplicities and discharge consumption) is deferred to the M3 runtime
layer where \texttt{Discharged} evidence is concrete.

\subsection{Summary of mechanization status}

\begin{table}[h]
\centering
\begin{tabular}{lll}
\toprule
Theorem & Status & File \\
\midrule
T1 (propositional, soundness) & \texttt{stated} (Rust impl. runs) & \texttt{Decidability.lean} \\
T2 (full) & \texttt{stub} & \texttt{Correspondence.lean} \\
T2 (symbolic half) & implicit in \texttt{Tamarin+ProVerif} & \\
T3 (full) & \texttt{stated} & \texttt{NonInterference.lean} \\
T4 (non-introduction) & \texttt{proven\_partial} & \texttt{ObligationSoundness.lean} \\
Substitution composition & \texttt{stated}; 4 structural lemmas proven & \texttt{Subst.lean} \\
Subject reduction & \texttt{stated} & \texttt{Reduce.lean} \\
Graded comonad laws & \texttt{proven} (3 laws) & \texttt{Graded.lean} \\
Indexed monad laws & \texttt{proven} (4 laws incl.\ strength) & \texttt{IndexedMonad.lean} \\
\bottomrule
\end{tabular}
\caption{Mechanization status as of May 2026.  Twelve theorems are
fully proven with no \texttt{sorry} and no new axioms.  Four headline
theorems are stated with proof closures detailed in
Section~\ref{sec:conclusion}.}
\label{tab:status}
\end{table}


\section{Symbolic Verification}
\label{sec:symbolic}
% Section 4 -- Symbolic Verification (Tamarin + ProVerif)
% Target: 3 pages.

The DLC wire protocol is modeled symbolically in two independent
provers: Tamarin (multiset rewriting + DPLL backward search) and
ProVerif (applied-pi + Horn-clause resolution).  Both verify the
load-bearing \emph{NonSplicing} property.

\subsection{The Tamarin model}

\texttt{models/tamarin/dlc.spthy} encodes five multiset-rewriting
rules:

\begin{itemize}
\item \texttt{Register\_pk} -- assigns each principal a long-term Ed25519 keypair.
\item \texttt{Reveal\_ltk} -- models static long-term-key compromise (Dolev-Yao+).
\item \texttt{Says} -- a principal signs a fresh proposition.
\item \texttt{SpeaksFor\_Issue} -- a principal signs ``$Q$ speaks for $P$''.
  Crucially, this rule emits \emph{both} a \texttt{SpeaksForIssued($P$,$Q$)}
  event and a \texttt{Says($P$, $\langle$speaks\_for, $P$, $Q\rangle$)} event,
  reflecting that a SpeaksFor token is just a particular shape of
  Says token.  Omitting the second event admits a self-delegation
  splice (see below).
\item \texttt{Delegate\_Accept} -- the verifier rule.  Its input patterns
  bind the same \texttt{\$Q} in the SpeaksFor token and the Says token,
  structurally enforcing the no-chain-splicing condition.
\end{itemize}

\paragraph{Lemmas verified.}  Three:

\begin{itemize}
\item \texttt{secrecy\_ltk}: the adversary never learns an honest
  principal's secret key.  \emph{Verified, 2 steps.}
\item \texttt{non\_splicing}: every \texttt{DelegateAccept($P$, $Q$, $\mathit{prop}$)}
  trace contains both a \texttt{SpeaksForIssued($P$, $Q$)} event and a
  \texttt{Says($Q$, $\mathit{prop}$)} event, unless $P$'s or $Q$'s key was
  revealed.  \emph{Verified, 18 steps.}  This is the property that
  defeats RFC 8693 chain-splicing.
\item \texttt{exec\_delegation}: a trace exists in which delegation
  is accepted without key reveal (executability sanity check).
  \emph{Verified, 7 steps.}
\end{itemize}

\paragraph{The self-delegation splice surfaced during modeling.}
On the first CI run with \texttt{SpeaksFor\_Issue} emitting only
\texttt{SpeaksForIssued} (not also \texttt{Says}), Tamarin found a
7-step counterexample to non-splicing.  The attack: take a SpeaksFor
token issued by $P$ naming some $Z$, repackage it as the Says half
of a delegation with $Q = P$ and $\mathit{prop} = \langle$speaks\_for,
$P$, $Z\rangle$.  Both halves verify (the signature was honestly
issued); the lemma's antecedent \texttt{Says($P$, $\mathit{prop}$)} is
not in the trace because \texttt{SpeaksFor\_Issue} only emitted
\texttt{SpeaksForIssued}.  Fix: emit both events.  This is the right
modeling because SpeaksFor IS a Says with a specific proposition shape.

\subsection{The ProVerif cross-check}

\texttt{models/proverif/dlc.pv} mirrors the Tamarin model in
applied-pi calculus.  Same three properties; different proof engine.
Two modeling gaps surfaced during cross-checking that the Tamarin
persistent-fact semantics had implicitly closed:

\begin{enumerate}
\item \emph{Secrecy was vacuously true} because the attacker can
  freely generate its own fresh \texttt{skey} values without going
  through \texttt{Register\_pk}.  Fix: added an \texttt{event
  HonestSk($sk$)} restricted to honestly-created keys; the secrecy
  query antecedent now restricts the universal quantifier to those
  keys.
\item \emph{NonSplicing was falsified} because the verifier accepted
  any \texttt{pkey} as $Q$, including attacker-fabricated ones.
  Fix: added a \texttt{table honest\_keys(pkey)} populated only
  during honest registration; the verifier process gates both
  halves of the chain with \texttt{get honest\_keys(=pkP) in;
  get honest\_keys(=pkQ) in}.
\end{enumerate}

Both fixes \emph{tightened the model to match DLC's actual wire
semantics} (the verifier checks against a pinned keyring; cf.\
\texttt{spec/threat-model.md} \S 3), not weakenings of the threat
model.  After fixes, ProVerif reports all three queries as
\texttt{is true} or the expected existence-style \texttt{is false}.

\subsection{Why two provers}

The Tamarin and ProVerif cross-check is the artifact's primary
methodological contribution to the symbolic-verification line.
Existing delegation-protocol work (e.g., AITH~\cite{chen2026aith})
verifies in Tamarin alone.  Different provers use different
intermediate representations and proof engines; a model bug that
happens to be invisible to one tool's heuristics typically surfaces
in the other.  In our case, the two gaps above were precisely such
bugs: ProVerif's applied-pi semantics surfaced what Tamarin's
\texttt{!Pk} persistent fact had implicitly assumed.

Whether to ship a result through two provers is a design choice
the field should normalize.  The cost---a second model file, ~150
LOC of applied-pi---is low.  The benefit---two-fold confidence and
explicit modeling-gap detection---is high.


\section{Wire Format and Runtime}
\label{sec:wire}
% Section 5 -- Wire Format and Runtime
% Target: 1 page.

DLC's wire format is a CBOR-encoded tagged-array tree.  Each \texttt{Term}
constructor maps to \texttt{[tag, arg1, arg2, \ldots]} where \texttt{tag}
is the constructor's stable index (0--21) and the arguments recurse.
Sub-types (\texttt{Prop}, \texttt{Principal}, \texttt{Label},
\texttt{Obligation}, \texttt{TimeBound}, \texttt{Signature}) use the
same tagged-array shape.

\paragraph{Content addressing.}  The canonical CBOR bytes of a sub-term
hash to a stable identifier (SHA-256, 32 bytes truncated to 8 in
display contexts).  A delegation chain is a Merkle DAG whose leaves
are signatures and whose internal nodes are \texttt{Delegate} and
\texttt{Attenuate} compositions.  Changing one leaf changes only the
hashes on the path to the root, enabling selective re-verification
and cross-organizational caching.  This is the same data-flow
property that makes macaroons and biscuits cacheable; DLC adds the
typed structure.

\paragraph{COSE\_Sign1 envelope.}  At the outermost layer, a DLC
token is a COSE\_Sign1 (RFC 9052) structure with the canonical
bytes as payload and an Ed25519 signature in the unprotected header.
The COSE infrastructure gives DLC tokens a JWT/CWT-style transport
shape for interoperability with existing token-handling systems
(OAuth, OIDC, Verifiable Credentials).

\paragraph{Implementation.}  The encoder and decoder are hand-rolled
against \texttt{ciborium::Value} rather than via Serde derive
macros, because the \texttt{dlc-core} crate must stay zero-dep to
remain Aeneas-translatable.  Encoder and decoder are exhaustive
pattern-matches against the term grammar; the round-trip property
$\mathsf{decode}(\mathsf{encode}(t)) = \mathsf{Some}(t)$ is verified by
22 unit tests covering every constructor plus compound nesting and
garbage rejection.

\paragraph{Tamarin/ProVerif exporters.}  The same runtime emits
Tamarin and ProVerif fact dumps via
\texttt{export\_tamarin::term\_to\_tamarin} and
\texttt{export\_proverif::term\_to\_proverif}.  A real runtime
witnessing a DLC term can thus emit a trace that the symbolic
provers' models could ingest---closing the operational half of L2.5
(wire $\leftrightarrow$ symbolic-term encoding).


\section{Computational Bridge}
\label{sec:computational}
% Section 6 -- Computational Bridge
% Target: 1 page.

DLC's symbolic non-splicing property is verified twice
(Section~\ref{sec:symbolic}); the computational reduction is the
load-bearing link to real cryptographic assumptions.

\subsection{The EUF-CMA reduction}

\begin{theorem}[T2 -- computational direction, sketch]
Let $A$ be a PPT adversary against DLC's symbolic non-forgery game
(\texttt{SymbolicSayForgery} in \texttt{models/easycrypt/Game.eca}).
There exists a PPT reduction adversary $A^*$ against the standard
EUF-CMA game (\texttt{EUF\_CMA}) such that:
\[
\Pr[\text{SymbolicSayForgery}(A) = 1] \leq \Pr[\text{EUF\_CMA}(A^*) = 1]
\]
By the EUF-CMA security of Ed25519, the right-hand side is
negligible, hence so is the left.
\end{theorem}

\paragraph{Game-hop sequence (sketched).}
\begin{enumerate}
\item $G_0$: \texttt{SymbolicSayForgery}.  Adversary $A$ given access to
  a signing oracle for a fresh keypair must produce a $(m^*, \sigma^*)$
  pair with $m^*$ not in the queried set and $\sigma^*$ verifying.
\item $G_1$: lazy-sample the signing oracle into a table.  Standard
  identical-until-bad rewrite.
\item $G_2$: replace the signing key with an EUF-CMA challenger.
  The reduction $A^*$ simulates $A$'s oracle by forwarding to its own
  EUF-CMA oracle.  $A$'s eventual forgery is also an EUF-CMA
  forgery.
\item Conclude: $\Pr[G_0(A)] = \Pr[G_2(A^*)] \leq \mathsf{Adv}^{\mathsf{EUF-CMA}}_{\mathsf{Ed25519}}(A^*)$.
\end{enumerate}

\paragraph{Mechanization path.}
Two viable routes:

\begin{enumerate}
\item \textbf{EasyCrypt.}  The skeleton in
  \texttt{models/easycrypt/Game.eca} is structurally complete: type
  declarations, signing oracles, the EUF-CMA module, the
  \texttt{SymbolicSayForgery} module, the EUF-CMA assumption as an
  axiom, and the reduction lemma documented (proof body absent;
  CI parse-checks via \texttt{easycrypt -batch}).  Closure
  requires cryptographer pairing (Blanchet was the named
  collaborator in our plan).

\item \textbf{VCVio.}  Tuma et al.~(2026)~\cite{tuma2026vcvio}
  mechanize EUF-CMA for Schnorr signatures in Lean 4 using a
  monadic oracle-effects framework.  Their proof template
  ports directly to Ed25519 (or to an abstract EUF-CMA-secure
  signature scheme).  This route keeps T2 inside the Lean
  development that holds T1, T3, and T4.
\end{enumerate}

Both routes deliver the same theorem.  VCVio is the autonomously
preferred route as of May 2026 because it stays in one proof
assistant; EasyCrypt is the publication-standard route because of
its cryptographer community.

\paragraph{Status.} The EasyCrypt skeleton parses against the
current grammar (CI step \texttt{easycrypt -batch Game.eca} passes,
6m36s).  The full reduction proof is open work.


\section{Related Work}
\label{sec:related}
% Section 7 -- Related Work
% Target: 2 pages.

DLC sits at the intersection of three research lineages: authorization
logic, capability-based access control, and AI agent delegation.

\subsection{Authorization logics with \texttt{says}-modality}

Garg \& Pfenning's constructive authorization logic~\cite{garg2006csf}
provides DLC's intuitionistic foundation: \texttt{says} as a modality
in a sequent calculus with cut elimination and non-interference proven
on paper.  Nexus Authorization Logic (NAL)~\cite{schneider2011nal}
mechanizes a related says-modality system in $\sim$3000 LOC of Coq.
DLC extends this lineage with explicit cryptographic carrier
semantics (\textsc{says-I} carries the signature intrinsically),
linear obligations ($\Box_O$), time modality ($\Diamond_\tau$), and
IFC labels.  Where NAL proves soundness against a belief-style
semantics, DLC's T2 targets a cryptographic correspondence---
arguably a stronger semantic anchor.

\subsection{Proof-carrying authorization}

Appel \& Felten~\cite{appel1999pcfs} introduce the
\emph{proof-carrying authorization} (PCA) paradigm: each request is
accompanied by a logical proof of its authorization.  PCFS, Bauer's
PCA-on-the-web work, and Garg's proof-carrying file system extend
this.  All use higher-order logic that is not decidable.  DLC's
contribution is occupying the decidability sweet-spot: the
propositional fragment admits a mechanized structural checker (T1),
and the design targets an $O(|M| \cdot \log |\Gamma|)$ bound for the
full calculus --- a target that is stated, not proven, at this
revision (Section~\ref{sec:metatheory}).

\subsection{Authenticated delegation for AI agents}

The agent-economy line begins with South et al.~(ICML
2025)~\cite{south2025authdelegation}: authenticated delegation as
OAuth+OIDC extensions.  AITH (Chen,
2026)~\cite{chen2026aith} mechanizes five Tamarin theorems
including post-quantum security.  LLMbda Calculus (Garby et al.,
2026)~\cite{garby2026llmbda} provides termination-insensitive
non-interference for an LLM-invoking lambda calculus.  DLC's
positioning: AITH covers the symbolic-protocol axis; LLMbda covers
the type-theoretic axis; DLC unifies both, plus the wire format and
cryptographic correspondence, into a single artifact with the
same chain-splicing-defeating rule appearing at each layer.

\subsection{Capability tokens}

Macaroons~\cite{birgisson2014macaroons} and
biscuits~\cite{couprie2022biscuits} provide cryptographically
verifiable, attenuable capability tokens.  ZCAP-LD~\cite{w3c2021zcap}
and UCAN~\cite{ucan2024} extend the capability-token model to
linked-data contexts.  These systems excel at wire-format design;
they lack a typed logic of obligations or non-interference.  DLC's
\texttt{attenuate} rule with the IFC-lattice side condition is the
natural lift of macaroon caveats to a typed setting.

\subsection{OAuth and the RFC 8693 lineage}

RFC 8693 (Token Exchange)~\cite{rfc8693} introduces the nested
\texttt{act} claim for delegation chains.  The February--March 2026 IETF
OAuth mailing-list thread on chain splicing~\cite{ietf2026splicing}
documented the cross-validation gap: an STS validating
\texttt{subject\_token} and \texttt{actor\_token} independently can
issue tokens asserting chains that never occurred.
draft-ietf-oauth-attenuating-agent-tokens-00~\cite{niyikiza2026attn}
proposes mitigations.  DLC's \texttt{delegate} rule is the
structural answer the IETF mitigation patterns implement procedurally:
the same-principal condition is a type-system invariant, not a
runtime check.

\subsection{Mechanized cryptography in Lean}

VCVio (Tuma et al.,
2026)~\cite{tuma2026vcvio} mechanizes EUF-CMA security for Schnorr
signatures in Lean 4 using monadic oracle effects.  This is the
closest Lean-side precedent for T2's computational reduction.  Earlier
Lean cryptographic work (\texttt{SymbolicCryptographyLean}) and
the Mathlib4 categorical infrastructure provide the substrate.

\subsection{Symbolic protocol verification with multiple provers}

Cremers \& Sasse~\cite{cremers2013tamarin} provide Tamarin.
Blanchet~\cite{blanchet2022proverif} provides ProVerif.  Joint use
of both has been done case-by-case (e.g., for TLS 1.3), but the
\emph{methodological argument} that running new protocols through
two provers detects model-overfitting bugs is, to our knowledge, not
made systematically in the delegation literature.  This paper
contributes one such systematic case: the two-modeling-gaps story
in Section~\ref{sec:symbolic}.


\section{Conclusion and Future Work}
\label{sec:conclusion}
% Section 8 -- Conclusion and Future Work
% Target: 1 page.

DLC unifies authorization logic, symbolic cryptography, computational
cryptography, runtime wire format, and reproducible CI in a single
artifact where the same chain-splicing-defeating rule appears at every
layer.  The unification is the contribution; the individual pieces
build on Garg-Pfenning (logic), Cremers-Blanchet (symbolic crypto),
and Tuma et al.~(VCVio, computational crypto in Lean).

\paragraph{What is proven.}  Twelve theorems with no \texttt{sorry}
and no new axioms:
\begin{itemize}
\item T4 non-introduction direction (\texttt{pendingObligations\_shift}
  closed; substitution-subset and step-case-analysis open).
\item Three graded comonad laws (identity, associativity, DpBudget
  saturating-add associativity).
\item Four strong indexed monad laws (left-id, right-id,
  associativity, strength left-unit).
\item Four structural substitution lemmas (\texttt{shift\_zero} plus three
  \texttt{substAt\_var\_*} cases).
\item Three Tamarin lemmas (secrecy, non-splicing, executability).
\item Three ProVerif queries (same, by independent prover).
\end{itemize}

\paragraph{What is stated but open.}  T1 full-calculus decidability,
T2 cryptographic correspondence, T3 non-interference, T4 multiset
form, substitution composition, subject reduction, the \S 4.4
protocol-logic correspondence.  Each has a documented closure path:
T1 via porting the Rust \texttt{decide\_pure} to Lean; T2 via VCVio
EUF-CMA reduction; T3 via hand-rolled logical relations in Lean
without Iris; T4 multiset form via the M3 runtime layer's
\texttt{Discharged} evidence; substitution composition via Ramos
et al.'s 709-line Metatheory framework once that becomes available.

\paragraph{What is not in this paper.}  We make no claim to having
established DLC as a paradigm.  That bar requires (a) full closure of
the open theorems, (b) external citation by adjacent work,
(c) standards-track adoption, (d) deployment beyond the reference
implementation.  Each of these depends on actors beyond the authors
and on multi-quarter time horizons.  We position the present paper as
a \emph{credible foundation} for that future work, with the artifact's
double-prover symbolic verification and reproducible ledger as
methodological contributions on their own.

\paragraph{Future work.}  In the medium term, four closure milestones
in priority order:

\begin{enumerate}
\item Close T4 non-introduction in full (substitution-subset lemma,
  step-case analysis).  Autonomously reachable in $\sim$2 weeks.
\item Close T1 propositional soundness via a Lean \texttt{decideLean}
  mirroring the Rust implementation.  Autonomously reachable in
  $\sim$3 weeks.
\item Close T2 via VCVio.  Autonomously reachable in $\sim$5-8 weeks.
\item Close T3 atomic-fragment + \texttt{lift} cases as a logical
  relations skeleton.  Autonomously reachable in $\sim$5 weeks;
  full T3 needs $\sim$3-5 months including the Lean infrastructure
  layer.
\end{enumerate}

In the longer term, collaborator engagement: Pfenning/Garg on T1/T3
proof technique, Cremers on Tamarin/ProVerif review, Blanchet on the
VCVio EUF-CMA reduction's tightness, Myers on IFC and obligation
algebra.  An IETF SECDISPATCH presentation of the wire-format draft
is the natural standards-track entry point.

\paragraph{The diagonal as the contribution.}  We close by stating
the contribution claim in its strongest form: DLC is the first
delegation calculus in which the same chain-splicing-defeating rule
appears across (a) a Lean-mechanized type system, (b) a Tamarin
symbolic model, (c) a ProVerif independent symbolic model, (d) an
EasyCrypt computational skeleton, (e) a Rust runtime with content-
addressed wire format.  Each layer is independently verifiable; their
agreement is the artifact's load-bearing methodological claim.  Whether
DLC ultimately earns the name of a paradigm depends on closures and
adoption that lie beyond this paper's horizon.  What this paper
contributes is the diagonal itself, with enough mechanical evidence
along each axis to make the unification credibly defensible under
external review.


\bibliographystyle{ACM-Reference-Format}
\bibliography{bib}

\end{document}
